\scanpage{229}\mspage{84}
Quand j'ordonne \deleted{$\unsure{P=\mathfrak{P}(W)}$} $\unsure{P}, Q\subset\unsure{\mathfrak{P}|\mathcal{P}}(\operatorname{Fl}(M))$
par la relation d'ordre induite par l'inclusion
de parties de $\operatorname{Fl}(M)$, on trouve des \inserted{\textsuperscript{ss}/}ensembles
ordonnés de $\mathfrak{P}(\operatorname{Fl}(M))$, stables par réunions
quelconques, et aussi par intersections quelconques
avec la seule réserve que dans le cas de
$P$, la famille \deleted{\unsure{doit}} d'indices doit être $\unsure{\ne|=}\varnothing$
(puisqu'alors on trouve comme intersection
$\operatorname{Fl}(M)$ \deleted{tout} entier, qui n'est pas en général
\unsure{contenue|contenu} dans $W$).

On va vérifier la condition $(\unsure{D_0|D})$ de la
page 69, pour $\mathcal{R}$ et les relations d'ordre.
Il y a deux vérifications à faire

$1^\circ)$ Donnons nous $A\in P$\inserted{, i.e. $A\subset W$}, et considérons l'ensemble
$\mathcal R(A)$ des $L\in Q$ tels que l'on ait \deleted{$\unsure{R\in}$}
$(A,L)\in\mathcal R$. Il est clair que c'est un ouvert
(et de $=$, que $\,{}^{\unsure{c|t}}\mathcal R(L)$ est un ouvert \inserted{de $\unsure{P|Q}$} pour
tout $L$ dans $Q$) -- i.e. $L\in R(A)$, $L'\subset L$
implique $L'\in R(A)$. Reste à trouver un
plus grand élément de $\mathcal R(A)$. Notons que
si $L\in R(A)$, alors \deleted{le chgt de base}\inserted{pour tout}
$p\in L$, $p:Y'\unsure{-|\to}Y$, le changt de base
$p^{*}:\operatorname{Cat}_{/Y}\unsure{-|\to}\operatorname{Cat}_{/Y'}\ :\quad X\unsure{\vert\!-\!\vert|\mapsto|\leftarrow}X\times_Y Y'$
doit transformer flèches $f\in A$ en flèches $g\in W$,
\diagram{\begin{tikzcd}[ampersand replacement=\&,column sep=small,row sep=small]
X_1\arrow[d,"f"'] \& X'_1\arrow[l,"\unsure{q_1|p_1}"']\arrow[d,"g"]\\
X_2\arrow[d] \& X'_2\arrow[l,"\unsure{q_2|p_2}"']\arrow[d]\\
Y \& Y'\arrow[l,"p"']
\end{tikzcd}}
comme il résulte du fait que
$L$ est stable par chgt de base.
Donc $L\subset\unsure{\mathbb{L}|L}_W(A)$, où on pose
% archive sheet number: 229
